Start with \wikilink{s27--b003}{maximal ideals and local rings}.
The following assertions concern commutative unital rings.

\begin{lemma}\label{maximal-ideal-existence}
Every proper ideal of a nonzero ring is contained in a maximal ideal.
In particular, every nonunit belongs to a maximal ideal.
\end{lemma}
\begin{proof}\label{maximal-ideal-existence-proof}
A chain of proper ideals containing the given ideal has proper union:
if $1$ belonged to the union, it would belong to a chain member.
Zorn's lemma gives a maximal member. For a nonunit $a$, apply this to
the proper ideal $(a)$.
\end{proof}
\begin{definition}\label{jacobson-radical}
The \emph{Jacobson radical} is
$J(R)=\bigcap_{\mathfrak m\text{ maximal}}\mathfrak m$.
For the zero ring the empty intersection is all of $R$.
\end{definition}
\begin{proposition}\label{jacobson-units}
If $a\in J(R)$, then $1-a$ is a unit. If $\mathfrak a\subseteq J(R)$,
every element of $1+\mathfrak a$ is a unit. Moreover,
$a\in J(R)$ exactly when $1-ra$ is a unit for every $r\in R$.
\end{proposition}
\begin{proof}\label{jacobson-units-proof}
A maximal ideal containing a nonunit $1-a$ would also contain $a$,
and hence $1$. Applying this observation to $-b$ and to $ra$ proves
the other forward assertions. Conversely, if $a\notin\mathfrak m$
for a maximal ideal, choose $r$ lifting the inverse of its residue
in $R/\mathfrak m$. Then $1-ra\in\mathfrak m$ is not a unit.
The zero ring satisfies the assertions directly.
\end{proof}
\begin{corollary}\label{local-jacobson}
If $(R,\mathfrak m)$ is local, then $J(R)=\mathfrak m$ and
$R^\times=R\setminus\mathfrak m$.
\end{corollary}
\begin{proof}\label{local-jacobson-proof}
The intersection has only one maximal ideal. Every nonunit is contained
in a maximal ideal, whereas no element of a proper ideal can be a unit.
\end{proof}
\begin{remark}\label{jacobson-nakayama-link}
Invertibility of $1+\mathfrak a$ is the ingredient in
\wikilink{a2s19--ed-general-nakayama--thm-general-nakayama}{Nakayama's lemma
over the Jacobson radical}.
\end{remark}
